\subsection*{33. A Concrete Positioning Worksheet}

For 1 buyer situation, complete the following analysis:

\begin{table}[htbp]
\centering
\small
\begin{tabularx}{\columnwidth}{>{\raggedright\arraybackslash}p{3.2cm} >{\raggedright\arraybackslash}X}
\toprule
\textbf{Question} & \textbf{Required Answer} \\
\midrule
Target buyer and trigger & Customer in a concrete event, not ``all professionals'' \\
Current alternative & Specific default workaround and commercial competitors \\
Inadequacy root cause & Observable error, latency, burden, or missed outcome \\
Distinctive capability & A real, deliverable production or teaching capability \\
Resulting customer value & Buyer-relevant outcome, ideally measurable \\
High-value segment & Circumstances where the difference shifts choice \\
Category frame & Truthful, familiar-enough market comparison frame \\
Available proof today & Inspectable demonstration, prior work, pilot evidence \\
Deliberate exclusions & Operational boundaries and accepted trade-offs \\
Economic deliverability & Capacity, contribution margin, and labor sustainability \\
\bottomrule
\end{tabularx}
\caption{The 10-point positioning specification worksheet.}
\label{tab:ch8-positioning-worksheet}
\end{table}

Produce several coherent candidate positions; do not assume the first is correct. Read the statements to real buyers and ask what they think the product does, what they would compare it with, and what evidence they would demand. If they misunderstand the category or invent benefits you cannot deliver, the frame needs revision.

\subsection*{34. Positioning Is an Internal Operating Contract}

A strong position constrains decisions. If a company claims to be an expert-led, review-intensive program, it must hire competent reviewers, maintain reasonable ratios, allocate review time, publish clear standards, and refuse enrollments it cannot support. If a product claims to be low-maintenance, engineering resources must prioritize onboarding reliability and supportability rather than demo-only features.

This is why a positioning exercise is insufficient when performed only by marketing staff. Product, delivery, sales, finance, and leadership must agree on what the claim obligates them to do. \textbf{A position that management will not finance is fiction. A position that salespeople must exaggerate is unstable. A position that operations cannot reproduce is a refund liability waiting to happen.}

\subsection*{35. The 5-Level Positioning Hierarchy}

For a multi-offering organization, distinguish 5 distinct levels:

\begin{itemize}
    \item \textbf{Institutional identity:} Why the organization exists and what enduring standards it aspires to embody.
    \item \textbf{Corporate position:} The basis on which the company competes and earns trust across all activities.
    \item \textbf{Portfolio architecture:} Relationships among different product families and divisions.
    \item \textbf{Product position:} Answers the specific buyer's choice question for a defined offering.
    \item \textbf{Campaign message:} Expresses 1 portion of that position for a specific channel or enrollment window.
\end{itemize}

A mission such as ``create the conditions in which human possibility becomes human capability'' does not tell an engineer whether to enroll in a course. A product message such as ``review your production AI architecture'' does not by itself communicate the ambition to become an enduring STEM research institution. Both can coexist when arranged at the correct level.

Organizations become confused when an aspirational institutional identity is used as if it were an immediately purchasable product, or when an entry-level product category becomes the company's permanent public definition.

\subsection*{36. Positioning Across Multiple Stakeholders}

An enterprise sale involves users, budget owners, security evaluators, finance, procurement, and managers whose careers are exposed to the decision. A program can be differentiated for engineers yet fail to persuade the executive responsible for paying for it. A college partnership can be attractive to a dean yet create unexpected burdens for the faculty members expected to deliver it.

The right approach is not to create 1 vague promise that pleases everyone. Define an economic or mission-relevant champion, then document specific value and risk for each stakeholder. For example, an AI education partner might promise the academic dean structured instructional delivery, the department faculty a manageable teaching workflow, the university administration auditable assessments, and the students authentic project practice. These are connected outcomes but not interchangeable claims.

The B2B position must make the purchasing decision \textbf{defensible}. Clear scope, documented deliverables, privacy and security boundaries, support commitments, and objective measurement may be more differentiating than extra curriculum topics.

\subsection*{37. Developing Positioning with Multiple Plausible Answers}

A founder can run a structured positioning sequence:

\begin{enumerate}
    \item Gather buyer evidence and actual alternatives from Chapters 6 and 7;
    \item Identify the 2 or 3 most consequential failures of the default alternative;
    \item List capabilities you have demonstrated or can reasonably build;
    \item Connect each capability to a buyer outcome and identify the strongest proof;
    \item Select a narrow group that disproportionately values that outcome;
    \item Generate 3 truthful category frames;
    \item Test unprompted comprehension and alternative recall with qualified buyers;
    \item Refine the offer and delivery process to make the best position true;
    \item Run priced decisions, not merely message-preference surveys; and
    \item Codify the position only after observing a repeated purchasing pattern.
\end{enumerate}

An output like ``for financial analysts who need auditable forecasting, an expert-reviewed modeling studio'' is a hypothesis. Its worth is determined by purchases, outcomes, and economics. A polished template is not evidence.

\subsection*{38. Avoid the False Dichotomy Between Positioning and Differentiation}

Some advisers argue that good positioning is all that matters; others insist the product must be materially better before messaging matters. The dichotomy is false. Positioning changes which benefits become salient in a comparison, while actual differentiation gives positioning something true to reveal. Sometimes changing the frame unlocks demand without altering the product. Sometimes the buyer's objection exposes a capability the company must build first.

Use positioning research as a diagnostic of \emph{what is already valuable and unrecognized}, and product development as a means of making the proposed position more valuable and durable. The 2 must coevolve under customer evidence.
